\documentclass[journal,onecolumn,11pt]{IEEEtran}
% NOTE: IEEE-стилско форматирање (IEEEtran). Компилирано со XeLaTeX поради
% поддршка на кирилица; распоредот/форматирањето е идентично со оригиналот,
% само фонтот е заменет со FreeSerif (Times-сличен, со полна кирилична поддршка).

\usepackage{fontspec}
\setmainfont{FreeSerif}
\setsansfont{FreeSans}
\setmonofont{FreeMono}
\usepackage{amsmath,amssymb}
\usepackage{babel}
\babelprovide[import, main]{macedonian}
\usepackage{siunitx}
\usepackage{booktabs}
\usepackage{array}
\usepackage{multirow}
\usepackage{graphicx}
\usepackage[siunitx]{circuitikz}
\usepackage{tikz}
\usetikzlibrary{arrows.meta,positioning,calc,shapes.geometric}
\usepackage[hidelinks,colorlinks=true,linkcolor=blue,citecolor=blue,urlcolor=blue]{hyperref}
\usepackage{url}
\usepackage{xcolor}
\usepackage{enumitem}
\usepackage{lettrine}
% IEEEtran's native \IEEEPARstart drop-cap uses TFM font loading that is
% incompatible with XeLaTeX/fontspec; redefine it with lettrine instead.
\renewcommand{\IEEEPARstart}[2]{\lettrine[lines=2,findent=2pt,nindent=0pt]{#1}{#2}}

\sisetup{detect-all,per-mode=symbol,group-digits=integer,range-phrase={ до }}
\DeclareSIUnit{\ohm}{\text{Ω}}

% Македонски имиња на структурните ознаки (IEEE-стил)
\AtBeginDocument{%
  \def\figurename{Сл.}%
  \def\tablename{ТАБЕЛА}%
  \def\refname{Референци}%
  \def\abstractname{Апстракт}%
  \def\IEEEkeywordsname{Клучни зборови}%
}
% потолерантно прекршување на редови за неделив кириличен текст
\tolerance=2000
\emergencystretch=2em
\hbadness=6000

% Convenience macros
\newcommand{\Vdcbus}{V_{\mathrm{bus}}}
\newcommand{\Rshunt}{R_{\mathrm{sh}}}
\newcommand{\Iphrms}{I_{\mathrm{ph}}}
\newcommand{\Fswsw}{f_{\mathrm{sw}}}

\begin{document}

\title{Дизајн на под-колата за мерење на фазна струја, мерење на напон на
DC-сабирницата и на фазите, и RC snubber за \mbox{20--\SI{60}{\volt}},
\mbox{3--\SI{5}{\kilo\watt} BLDC/PMSM} инвертор со векторско управување (FOC)}

\author{Извештај за хардверски дизајн\\
\small Целен контролер: STMicroelectronics STM32G473RCT6 \quad
Енергетски степен: Infineon IPT015N10N5 (\SI{100}{\volt}, \SI{1.5}{\milli\ohm}) +
Texas Instruments UCC27211A-Q1\\
\small Склопна фреквенција \SI{40}{\kilo\hertz} \quad Набавка на компоненти: LCSC Electronics}


\markboth{Извештај за дизајн на аналогниот влезен степен на BLDC/PMSM инвертор}{}
\maketitle

\begin{abstract}
Овој извештај ја презентира анализата, изведувањето и изборот на компоненти за три
аналогни под-кола од влезниот степен на инвертор за трифазен безчеткест еднонасочен /
синхрон мотор со перманентни магнети (BLDC/PMSM), напојуван од батерија: (i) ланецот за
мерење на струја заснован на in-line (in-phase) shunt, (ii) ланците за мерење на напон на
DC-сабирницата и на секоја фаза, и (iii) RC snubber на склопниот јазол. Системот е напојуван
или од \SI{24}{\volt} батерија (\mbox{20--\SI{30}{\volt}}, \SI{3}{\kilo\watt} мотор) или од
\SI{48}{\volt} батерија (\mbox{40--\SI{60}{\volt}}, \SI{5}{\kilo\watt} мотор) и работи на
\SI{40}{\kilo\hertz}. Континуираната фазна струја во најлош случај е изведена
(\(\approx\SI{136}{\ampere}\) RMS, \(\approx\SI{193}{\ampere}\) врв) и искористена за
димензионирање на shunt од метален елемент со \SI{0.1}{\milli\ohm}, читан од струен засилувач
Texas Instruments INA240A2 со подобрено потиснување на pulse-width-modulation (PWM). Споредени
се четирите топологии за мерење (low-side shunt, high-side shunt, in-line shunt и
магнетна/Hall), а in-line топологијата е оправдана за целосно векторско управување (FOC).
Прашањето дали напонските разделници бараат бафер со операционен засилувач (op-amp) е одговорено
квантитативно врз основа на моделот на семплирање на аналогно-дигиталниот конвертор (ADC) на
STM32G4, и е покажано дека е непотребен бидејќи изворната импеданција на разделникот
(\(\approx\SI{4.5}{\kilo\ohm}\)) се полни до 12-битна точност во рамки на време на семплирање
кое е тривијално достапно, додека шесте внатрешни op-amp засилувачи на MCU остануваат како
резервно решение без дополнителен трошок. Snubber-от е димензиониран од објавената излезна
капацитивност на уредот и од стандардниот метод со две мерења за екстракција на ѕвонењето
(ringing), и е поставен како do-not-populate (DNP) footprint. Сите вредности на компонентите
се пресметани, секоја избрана компонента е поврзана со LCSC нарачкин број, а сите инженерски
тврдења се поткрепени со цитати до примарни извори од производителите или признати референци со
активни линкови.
\end{abstract}

\begin{IEEEkeywords}
BLDC, PMSM, векторско управување (FOC), in-line мерење на струја, shunt отпорник, INA240,
RC snubber, напонски разделник, STM32G4, изворна импеданција на ADC, EMI филтрирање.
\end{IEEEkeywords}

\section{Вовед и опсег}
\IEEEPARstart{Ц}{елниот} систем е неизолиран трифазен напонски инвертор со заедничко
заземјување (common-ground), наменет да управува со кој било BLDC/PMSM мотор под сензорско
управување (Hall или квадратурен енкодер) или безсензорско (повратна електромоторна сила,
back-EMF) FOC. Контролерот е STM32G473RCT6 --- motor-control MCU базиран на
Arm\textsuperscript{\textregistered} Cortex\textsuperscript{\textregistered}-M4 кој интегрира
пет \mbox{12-битни} ADC-а, шест операционни засилувачи употребливи во режим со програмабилно
засилување (PGA) или како follower, седум аналогни компаратори, три напредни тајмери за
управување со мотор и бафер за напонска референца (VREFBUF)~\cite{stm32g473ds}. Gate driver-от
е Texas Instruments UCC27211A-Q1~\cite{ucc27211}, а секој прекинувач е Infineon
OptiMOS\textsuperscript{\texttrademark}~5 IPT015N10N5 (\SI{100}{\volt}, \SI{1.5}{\milli\ohm},
\SI{300}{\ampere})~\cite{ipt015ds,ipt015part}. Енергетскиот степен (gate driver-от и изборот на
MOSFET) и MCU се фиксни влезни услови за оваа работа.

Овој извештај ги дизајнира трите преостанати аналогни блока:
\begin{enumerate}[leftmargin=1.4em]
  \item ланецот за \emph{мерење на фазна струја} (in-line shunt + засилувач + ADC интерфејс), за
        сите три фази;
  \item ланците за \emph{мерење на напон} на DC-сабирницата и на секој од трите фазни јазли; и
  \item \emph{RC snubber} паралелно со енергетските прекинувачи.
\end{enumerate}
За секој блок извештајот (а) го објаснува принципот на работа и потребните компоненти,
(б) ги изведува управувачките равенки, (в) пресметува нумерички вредности за тековната
спецификација, (г) избира конкретна компонента достапна од LCSC Electronics, и (д) го оправдува
изборот. Две конкретни прашања поставени од спецификацијата се обработени детално: споредбата на
топологиите за мерење на струја што го мотивира in-line изборот (Дел~\ref{sec:csurvey}), и
потребата (или непотребата) од op-amp бафер пред ADC (Дел~\ref{sec:buffer}). Филтрирањето на
шумот е третирано таму каде што се генерира секој сигнал и е консолидирано во Дел~\ref{sec:noise}.

\subsection{Ограничување за набавка (sourcing)}
Секоја избрана компонента мора да биде нарачлива од LCSC Electronics, при што се претпочитаат
компоненти со голем обрт за да биде реална нарачка со количини од \(500+\). LCSC држи на залиха
повеќе од \num{560000} SKU од \(2600+\) производители~\cite{lcschome}; верифицираните активни
компоненти и стандардните пасивни компоненти употребени тука се производи со голем обем. Бидејќи
залихата кај дистрибутерот се менува секојдневно, спецификацијата на материјали во
Дел~\ref{sec:bom} ја означува единствената специјална компонента (shunt-от со голема моќност) за
потврда на залиха при пуштање на нарачката и нуди алтернатива составена од крајно вообичаени
компоненти.

\section{Работен опсег и фазна струја во најлош случај}\label{sec:env}
Хардверот за мерење на струја мора да биде димензиониран за најголемата фазна струја што
инверторот може да ја носи, што ги одредува вредноста на shunt-от и засилувањето на засилувачот.
Табела~\ref{tab:env} ги наведува четирите работни точки на напојување/моќност.

Електричната влезна моќност \(P\) е еднаква на излезната AC моќност на инверторот поделена со
ефикасноста на конверторот \(\eta\); земањето \(\eta\!\to\!1\) е конзервативната претпоставка (што
ја максимизира струјата). За синусоидно модулиран трифазен инвертор со просторно-векторска PWM
(SVPWM), максималната амплитуда на основната хармоника на фазниот напон е
\(\hat V_{\mathrm{ph}}=\Vdcbus/\sqrt{3}\), па максималниот RMS фазен напон е
\begin{equation}
V_{\mathrm{ph,rms}}=\frac{\Vdcbus}{\sqrt{6}} .
\label{eq:vph}
\end{equation}
Со три фази и фактор на моќност на поместување \(\mathrm{PF}=\cos\varphi\), AC моќноста е
\(P = 3\,V_{\mathrm{ph,rms}}\,\Iphrms\,\mathrm{PF}\). Со замена на \eqref{eq:vph} и решавање по
RMS фазната струја,
\begin{equation}
\boxed{\;\Iphrms = \frac{\sqrt{6}}{3}\,\frac{P}{\Vdcbus\,\mathrm{PF}}
        \approx 0.816\,\frac{P}{\Vdcbus\,\mathrm{PF}}\;}
\label{eq:iph}
\end{equation}
која е најголема при минимален напон на сабирницата и минимален фактор на моќност.
Равенката~\eqref{eq:iph} ја претставува струјата во номиналната работна точка; да се забележи дека
мотор при ниска брзина и висок вртежен момент може да бара споредлива фазна струја при
\emph{пониска} моќност на сабирницата, па \eqref{eq:iph} се користи како термичка/континуирана
проектна точка, а дополнителна резерва се остава за преодни и stall услови во опсегот на
засилувачот (Дел~\ref{sec:csa}).

\begin{table}[t]
\centering
\caption{Работни точки и резултантна фазна струја од \eqref{eq:iph} (\(\mathrm{PF}=0.9\)).}
\label{tab:env}
\small
\begin{tabular}{l c c c c}
\toprule
Напојување / мотор & \(\Vdcbus\) (V) & \(P\) (kW) & \(\Iphrms\) (A,rms) & \(\hat I_{\mathrm{ph}}\) (A,врв)\\
\midrule
\SI{24}{\volt} бат., \SI{3}{\kilo\watt} & 20 (мин) & 3 & \textbf{136.1} & \textbf{192.5}\\
\SI{24}{\volt} бат., \SI{3}{\kilo\watt} & 30 (макс) & 3 & 90.7 & 128.3\\
\SI{48}{\volt} бат., \SI{5}{\kilo\watt} & 40 (мин) & 5 & 113.4 & 160.4\\
\SI{48}{\volt} бат., \SI{5}{\kilo\watt} & 60 (макс) & 5 & 75.6 & 106.9\\
\bottomrule
\end{tabular}
\end{table}

Ограничувачката работна точка е \SI{3}{\kilo\watt} при \SI{20}{\volt}:
\(\Iphrms^{\max}\approx\SI{136}{\ampere}\) RMS, \(\hat I_{\mathrm{ph}}^{\max}\approx\SI{193}{\ampere}\)
врв. Ова е под номиналната граница од \SI{300}{\ampere} за пакувањето на IPT015N10N5~\cite{ipt015part},
што потврдува дека енергетскиот степен е во согласност со струјната цел. Сите вредности за мерење на
струја подолу се однесуваат на оваа работна точка.

\section{Мерење на фазна струја}\label{sec:current}

\subsection{Преглед и споредба на топологии за мерење на струја}\label{sec:csurvey}
Постојат четири фамилии на мерење на струја кај инвертор; нивниот мерен јазол е прикажан на
Сл.~\ref{fig:topo} за еден half-bridge.

\begin{figure}[t]
\centering
\begin{circuitikz}[scale=0.92,transform shape,american]
% bus rails
\draw (0,4) node[anchor=east]{$\Vdcbus$} -- (6.5,4);
\draw (0,-1.2) node[anchor=east]{GND} -- (6.5,-1.2);
% high-side switch (box) with optional high-side shunt above
\draw (2,3.0) rectangle (3,3.8) node[midway]{$Q_H$};
\draw (2.5,4) -- (2.5,3.8);
\draw (2.5,3.0) -- (2.5,1.7);
% low-side switch
\draw (2,0.9) rectangle (3,1.7) node[midway]{$Q_L$};
\draw (2.5,1.7) -- (2.5,1.7);
% phase node
\draw (2.5,1.7) node[circ]{} -- (4.6,1.7) node[anchor=west]{фаза $\to$ мотор};
% inline shunt in phase
\draw (3.6,1.7) node[anchor=south]{\scriptsize inline $\Rshunt$};
\draw[fill=yellow!30] (3.25,1.55) rectangle (3.95,1.85);
% low-side shunt to ground
\draw (2.5,0.9) -- (2.5,0.2);
\draw[fill=yellow!30] (2.15,-0.1) rectangle (2.85,0.2) node[midway]{};
\draw (2.5,0.0) node[anchor=west]{\scriptsize low-side $\Rshunt$};
\draw (2.5,-0.1) -- (2.5,-1.2);
% high-side shunt position marker
\draw[fill=yellow!30] (2.15,4.0) rectangle (2.85,4.3);
\draw (2.5,4.45) node[anchor=south]{\scriptsize high-side $\Rshunt$ (сабирница)};
\draw (2.5,4.3) -- (2.5,5.0) -- (5.5,5.0) -- (5.5,4) ;
% DC-link single shunt
\draw[fill=orange!25] (5.15,-1.55) rectangle (5.85,-1.25);
\draw (5.5,-1.7) node[anchor=north]{\scriptsize единечен DC-link $\Rshunt$};
\draw (5.5,-1.2)--(5.5,-1.25); \draw(5.5,-1.55)--(5.5,-2.0)--(0,-2.0);
\end{circuitikz}
\caption{Можни позиции на shunt во една гранка од инверторот: high-side (на сериска врска со
гранката од сабирницата), low-side (на сериска врска со гранката кон заземјувањето), in-line (на
сериска врска со фазниот спроводник), и единечниот DC-link shunt во заедничкиот повраток.
Hall/магнетната алтернатива го заменува shunt-от со сензор на поле околу фазниот спроводник.}
\label{fig:topo}
\end{figure}

\begin{table*}[t]
\centering
\caption{Споредба на топологии за мерење на струја кај трифазен инвертор.}
\label{tab:topo}
\small
\renewcommand{\arraystretch}{1.25}
\begin{tabular}{p{2.6cm} p{2.4cm} p{2.6cm} p{3.0cm} p{3.4cm}}
\toprule
\textbf{Атрибут} & \textbf{Low-side shunt} & \textbf{High-side shunt} & \textbf{In-line (in-phase) shunt} & \textbf{Магнетна / Hall} \\
\midrule
Common-mode напон на мерниот елемент & Близу до земја ($\approx\!0$) & На/близу $\Vdcbus$ & Се ниша од шина до шина на $\Fswsw$ со високо $\mathrm{d}v/\mathrm{d}t$ & Галвански изолиран; сензорот не гледа CM \\
Ја мери вистинската фазна струја при сите работни циклуси & Не --- само додека спроведува $Q_L$; слеп при $\approx100\%$ duty & Не --- само додека спроведува $Q_H$ & \textbf{Да --- континуирано, при кој било duty} & Да \\
Барање за засилувач & Low-CM CSA & High-CM CSA & High-CM CSA со PWM потиснување & --- (сензорот е засилувачот)\\
DC способност & Да & Да & Да & Да (Hall) / Не (CT, Rogowski)\\
Загуба од вметнување & $I^2R$ во shunt & $I^2R$ во shunt & $I^2R$ во shunt & Занемарлива\\
Точност / дрифт & Висока / низок & Висока / низок & Висока / низок & Умерена; офсет и термички дрифт, нелинеарност\\
Изолација & Не & Не & Не & Да\\
Релативна цена & Најниска & Ниска & Ниска--умерена (CSA/фаза) & Највисока\\
Најпогодна за & 6-step / евтини, ограничен duty опсег & Мониторинг на батерија/полнач & \textbf{Високи перформанси FOC низ целиот опсег на брзина/duty} & Многу високи струи или дизајни каде изолацијата е задолжителна\\
\bottomrule
\end{tabular}
\end{table*}

Решавачкиот атрибут за FOC е можноста да се семплира \emph{вистинската} фазна струја во секој PWM
период, независно од работниот циклус. Low-side shunt-от носи фазна струја само додека спроведува
low-side уредот, па при висок duty (висока брзина / висока модулација) прозорецот на спроведување
се намалува и мерењето станува непрецизно или бара изобличување на PWM (вметнување на минимално
време на спроведување)~\cite{ti_lowside_threephase}. Единечниот DC-link shunt ги реконструира трите
фазни струи од моменталната склопна состојба, но трпи од ``ненабљудливи'' области во близина на
границите на напонските вектори каде што времињата на задржување на два активни вектори паѓаат под
времето на стабилизирање/blanking на shunt-от, што повторно бара модификација на PWM. In-line
(in-phase) shunt-от е поставен на сериска врска со изводот на моторот, па ја спроведува фазната
струја континуирано и може да се семплира при кој било работен циклус; затоа токму оваа топологија
се користи во референтните дизајни за серво и FOC со високи перформанси~\cite{tida00913,ina240ds}.
Нејзината единствена тешкотија е тоа што мерниот елемент се наоѓа на јазол што се ниша помеѓу земја и
$\Vdcbus$ на $\Fswsw$ со големо $\mathrm{d}v/\mathrm{d}t$, па засилувачот мора да комбинира широк
common-mode (CM) опсег со силно потиснување на CM преодните појави~\cite{ina240ds,tida00913}.

Магнетната/Hall алтернатива ја отстранува загубата од вметнување и обезбедува изолација, но по
повисока цена и со поголем офсет и термички дрифт отколку ланецот shunt-плус-zero-drift-засилувач;
изолацијата не е потребна тука бидејќи дизајнот е експлицитно common-ground. Со оглед на
common-ground ограничувањето, потребата за FOC набљудливост при целосен duty на сите три фази, и
управливата дисипација во shunt-от изведена подолу, избрана е топологијата \textbf{in-line shunt},
во согласност со спецификацијата и со 48-V in-line FOC референтниот дизајн на TI,
TIDA-00913~\cite{tida00913}.

\subsection{Димензионирање на shunt отпорникот}\label{sec:shunt}
Три конкурентни ограничувања ја одредуваат вредноста на shunt-от $\Rshunt$: (1) дисипацијата на
моќност $P_{\Rshunt}=\Iphrms^2\Rshunt$ мора да остане во рамки на економичен површински-монтиран
елемент; (2) сигналниот напон на полна скала $\hat V_{\mathrm{sh}}=\hat I_{\mathrm{ph}}\Rshunt$ мора,
по засилувањето, да се вклопи во влезниот прозорец на ADC без отсекување (clipping) при врвот во
најлош случај (со резерва за преоптоварување); и (3) $\Rshunt$ треба да биде доволно голем за да
офсетот и шумот на засилувачот, сведени на струја, останат мали.

За $\Rshunt=\SI{0.1}{\milli\ohm}$ при работната точка во најлош случај
($\Iphrms^{\max}=\SI{136}{\ampere}$ rms):
\begin{align}
P_{\Rshunt} &= \Iphrms^{2}\Rshunt = (136.1)^2\times 0.1\times10^{-3} = \SI{1.85}{\watt},\\
\hat V_{\mathrm{sh}} &= \hat I_{\mathrm{ph}}\,\Rshunt = 192.5\times 0.1\times10^{-3} = \SI{19.3}{\milli\volt}.
\end{align}
Отпор од \SI{0.1}{\milli\ohm} ја држи дисипацијата на \SI{1.85}{\watt}, што удобно се справува од
страна на shunt од метален елемент (метална лента) со \SI{5}{\watt} (резерва од $2.7\times$), додека
\SI{19.3}{\milli\volt} врв е доволно голем за прецизно засилување. Елемент од метална лента е избран
поради неговата многу ниска паразитна индуктивност (ESL) и низок температурен коефициент;
четири-приклучна (Kelvin) врска е задолжителна за сигналниот напон да се зема само на резистивниот
елемент, а не на лемните споеви со голема струја~\cite{ina240ds}. ESL е важна бидејќи измерениот
напон е $v_{\mathrm{sh}}(t)=\Rshunt\,i(t)+\mathrm{ESL}\,\dfrac{\mathrm{d}i}{\mathrm{d}t}$; на
\SI{40}{\kilo\hertz} со брзи склопни рабови индуктивниот член мора да биде занемарлив наспроти
резистивниот член од \SI{0.1}{\milli\ohm}, што го постигнува елемент од метална лента со
$<\SI{2}{\nano\henry}$.

\subsection{Избор на струен засилувач и засилување}\label{sec:csa}
Засилувачот мора (i) да поднесува CM напон што се ниша од малку под земја (спроведување на body-диода
/ dead-time, $\approx\SI{-1}{\volt}$) до $\Vdcbus=\SI{60}{\volt}$ плюс пречекорувањето од ѕвонење;
(ii) да го потиснува големото CM $\mathrm{d}v/\mathrm{d}t$ при секој склопен раб така што излезот да
не ѕвони; (iii) да биде \emph{двонасочен} бидејќи фазната струја е AC; и (iv) да има пропусен опсег
далеку над \SI{40}{\kilo\hertz}. Texas Instruments INA240 ги задоволува сите четири: тоа е струен
засилувач со напонски излез и zero-drift архитектура, со CM опсег од $-4$ до $+\SI{80}{\volt}$
\emph{независно од напојувањето}, со ``enhanced PWM rejection'' експлицитно проектирано да ги потисне
CM преодните појави на PWM за погон на мотор, и се нуди во четири фиксни засилувања од
$20,50,100,200$~V/V~\cite{ina240ds}. Долната граница од $-\SI{4}{\volt}$ ја прифаќа екскурзијата под
земја, а горната граница од $+\SI{80}{\volt}$ остава \SI{20}{\volt} резерва над сабирницата од
\SI{60}{\volt} за пречекорувањето од ѕвонење.

При напојување од \SI{3.3}{\volt} аналоген извор, двонасочниот излез при нулта струја е центриран од
двата референтни приклучоци на INA240, чиј внатрешен разделник ја поставува средната точка на излезот
на $V_{\mathrm{ref}}=(V_{\mathrm{REF1}}+V_{\mathrm{REF2}})/2$~\cite{ina240ds}; врзувањето на
$\mathrm{REF1}$ за земја и $\mathrm{REF2}$ за шината од \SI{3.3}{\volt} дава
$V_{\mathrm{ref}}=\SI{1.65}{\volt}$. Тогаш излезниот напон е
\begin{equation}
V_{\mathrm{out}} = V_{\mathrm{ref}} + G\,\Rshunt\,i_{\mathrm{ph}} .
\label{eq:csaout}
\end{equation}
Со rail-to-rail излез, употребливата едностранска амплитуда е $\approx\SI{1.45}{\volt}$ (со резерва од
\SI{0.2}{\volt} до секоја шина). Со избор на засилување $G=50$ (INA240A2), линеарната врвна струја на
полна скала е
\begin{equation}
\hat I_{\mathrm{FS}} = \frac{1.45}{G\,\Rshunt}
   = \frac{1.45}{50\times 0.1\times10^{-3}} = \SI{290}{\ampere}\ \text{врв}\;(\SI{205}{\ampere}\,\text{rms}).
\end{equation}
Врвот во најлош случај од \SI{193}{\ampere} произведува излезна екскурзија од
$G\Rshunt\hat I_{\mathrm{ph}}=50\times0.1\text{m}\times192.5=\SI{0.96}{\volt}$, т.е.\
$V_{\mathrm{out}}\in[0.69,2.61]\,\si{\volt}$, оставајќи $\approx\SI{0.49}{\volt}$ пред отсекувањето.
Оваа резерва е намерно зачувана за преодно пречекорување и овозможува прагот за пречекорна струја
(over-current, OC) да се постави на $\approx\SI{250}{\ampere}$ врв ($1.3\times\hat I_{\mathrm{ph}}^{\max}$,
излезна екскурзија \SI{1.25}{\volt}), што може да се детектира од аналоген компаратор на STM32G4 или
софтверски пред засилувачот да засити. Квантизацијата и офсетот, сведени на струја, се
\begin{align}
\Delta I_{\mathrm{LSB}} &= \frac{V_{\mathrm{ref+}}/2^{N}}{G\,\Rshunt}
   = \frac{3.3/4096}{50\times0.1\text{m}} = \SI{0.16}{\ampere/LSB},\\
I_{\mathrm{os}} &= \frac{V_{\mathrm{os(max)}}}{\Rshunt} = \frac{\SI{25}{\micro\volt}}{0.1\text{m}\Omega}
   = \SI{0.25}{\ampere},
\end{align}
двете занемарливи наспроти сигнал од \SI{136}{\ampere} (zero-drift архитектурата на INA240 дава
$V_{\mathrm{os}}\le\SI{25}{\micro\volt}$~\cite{ina240ds}). Поради тоа, парот засилување-50 /
\SI{0.1}{\milli\ohm} се претпочита пред алтернативниот пар засилување-20 / \SI{0.25}{\milli\ohm}, кој
би дисипирал \SI{4.6}{\watt} по shunt за иста полна скала.

\subsection{ADC интерфејс и излезно филтрирање}\label{sec:csafilter}
За FOC фазната струја се семплира \emph{синхроно} со PWM носителот (вообичаено при врвот/долот на
бројачот) така што моментот на семплирање се совпаѓа со средината на склопната состојба, каде што
компонентата на бранувањето (ripple) на индукторот е на нејзината средна вредност; ова природно го
потиснува бранувањето на \SI{40}{\kilo\hertz} и неговите хармоници без тежок филтер~\cite{an5346}.
Тежок anti-alias филтер наместо тоа би додал фазно доцнење на токму оној сигнал на кој се затвора
струјната јамка, влошувајќи ја јамката. Затоа излезот на CSA користи само лесен RC филтер од прв ред
($R_f=\SI{100}{\ohm}$, $C_f=\SI{2.2}{\nano\farad}$) кој (i) го ограничува опсегот на високофреквентниот
склопен шум и ѕвонење (гранична фреквенција $f_c=1/2\pi R_fC_f=\SI{723}{\kilo\hertz}$, далеку над
содржината на струјната јамка во kHz-опсегот, но придушувајќи го ѕвонењето од десетици MHz), и (ii)
дејствува како резервоар на полнеж за sample-and-hold на ADC. Ниската излезна импеданција на INA240 ги
управува оваа мрежа и кондензаторот за семплирање на ADC со обилна резерва за стабилизирање.
Хардверското oversampling/усреднување на STM32G4 може да се овозможи за дополнително намалување на
шумот по цена на пропусност~\cite{an5346}. Сл.~\ref{fig:csa} го прикажува комплетниот ланец по фаза.

\begin{figure}[t]
\centering
\begin{circuitikz}[scale=1.0,transform shape,american]
% shunt row (top)
\draw (0,4.7) node[anchor=east]{фазен јазол} -- (1.3,4.7);
\draw[fill=yellow!30] (1.3,4.5) rectangle (2.5,4.9);
\draw (1.9,5.2) node{\scriptsize $R_{\mathrm{sh}}=\SI{0.1}{\milli\ohm}$ (Kelvin)};
\draw (2.5,4.7) -- (3.9,4.7) node[anchor=west]{кон мотор};
% Kelvin taps (non-crossing): right terminal -> IN+ (upper), left terminal -> IN- (lower)
\draw (2.3,4.5) -- (2.3,3.0) -- (2.7,3.0);
\draw (1.5,4.5) -- (1.5,2.5) -- (2.7,2.5);
\draw (2.45,3.0) node[anchor=south]{\scriptsize IN+};
\draw (2.0,2.5) node[anchor=south]{\scriptsize IN--};
% INA240 block
\draw (2.7,1.9) rectangle (5.1,3.5);
\draw (3.9,2.45) node{\small INA240A2};
\draw (3.9,3.05) node{\scriptsize $G=50$};
% supply stub (top, label kept clear of the shunt row)
\draw (3.4,3.5) -- (3.4,4.0) node[anchor=east]{\scriptsize $3V3_{\mathrm A}$+\SI{100}{\nano\farad}};
% reference stub (bottom)
\draw (3.9,1.9) -- (3.9,1.2) node[anchor=north,align=center]{\scriptsize REF1=GND, REF2=3V3\\\scriptsize $\Rightarrow V_{\mathrm{ref}}=\SI{1.65}{\volt}$};
% output + filter (right)
\draw (5.1,2.7) -- (5.6,2.7) to[R,l={$R_f$, \SI{100}{\ohm}}] (7.1,2.7);
\draw (7.1,2.7) node[circ]{} to[C,l={$C_f$, \SI{2.2}{\nano\farad}}] (7.1,1.5) node[ground]{};
\draw (7.1,2.7) -- (7.8,2.7) node[anchor=west]{\scriptsize ADC};
\end{circuitikz}
\caption{In-line ланец за мерење на струја по фаза: Kelvin-приклучен shunt од метална лента со
\SI{0.1}{\milli\ohm}, INA240A2 ($G=50$) со референца на средна скала, и лесен RC од
\SI{723}{\kilo\hertz} во ADC на STM32G4. Се користат три идентични ланци (фази U, V, W).}
\label{fig:csa}
\end{figure}

\section{Мерење на напон}\label{sec:voltage}

\subsection{Барања}
Контролерот мора да поддржува сензорско FOC (Hall/енкодер) и безсензорско FOC, при што второто ја
реконструира положбата на роторот од back-EMF. Затоа е потребно хардверско мерење на напон за (i)
напонот на DC-сабирницата --- за заштита од пренапон/поднапон, скалирање на индексот на модулација и
пресметка на моќност --- и (ii) секој од трите напони на фазните јазли --- за набљудување на back-EMF
и за повратна врска по фазен напон. Сите четири канали го делат заземјувањето на инверторот
(неизолиран дизајн), па секој е заземјен отпорнички разделник што напојува ADC канал на STM32G4.
Каналот на сабирницата е спор (заштита/скалирање), додека фазните канали мора да го зачуваат
тајмингот на back-EMF за набљудувачот на положбата.

\subsection{Дизајн на отпорнички разделник}\label{sec:divider}
Опсегот на конверзија на ADC е \(0\)--\(V_{\mathrm{REF+}}\). Користењето
\(V_{\mathrm{REF+}}=V_{\mathrm{DDA}}=\SI{3.3}{\volt}\) го максимизира опсегот; сооднос на разделникот
\(k=(R_{\mathrm{top}}+R_{\mathrm{bot}})/R_{\mathrm{bot}}\) го пресликува максималниот измерен напон во
близина на полна скала со резерва. Со избор на \(R_{\mathrm{top}}=\SI{110}{\kilo\ohm}\) и
\(R_{\mathrm{bot}}=\SI{4.7}{\kilo\ohm}\),
\begin{equation}
k=\frac{R_{\mathrm{top}}+R_{\mathrm{bot}}}{R_{\mathrm{bot}}}
 =\frac{110+4.7}{4.7}=24.4,\qquad V_{\mathrm{ADC}}=\frac{V_{\mathrm{in}}}{k}.
\end{equation}
Ова пресликува \(\Vdcbus=\SI{60}{\volt}\to\SI{2.46}{\volt}\) (\(\SI{73.2}{\volt}\) при полна скала,
т.е.\ \SI{22}{\percent} резерва над опсегот) и \(\SI{30}{\volt}\to\SI{1.23}{\volt}\). Истиот сооднос се
користи за фазните разделници бидејќи фазен јазол може да го достигне напонот на сабирницата.
Постојаната струја и дисипацијата при \SI{60}{\volt} се
\begin{equation}
I_{\mathrm{div}}=\frac{60}{R_{\mathrm{top}}+R_{\mathrm{bot}}}=\SI{523}{\micro\ampere},\quad
P_{R_{\mathrm{top}}}=I_{\mathrm{div}}^2R_{\mathrm{top}}=\SI{30}{\milli\watt},
\end{equation}
па сите четири разделници заедно повлекуваат \(\approx\SI{126}{\milli\watt}\) --- занемарливо. Горниот
отпорник го гледа целиот напон на сабирницата; единечен thick-film чип од \SI{110}{\kilo\ohm} со
номинал \(\ge\SI{150}{\volt}\) (на пр.\ 0805/1206) е адекватен при \SI{60}{\volt}, но се препорачува
делење на \(R_{\mathrm{top}}\) на два сериски делови од \SI{55}{\kilo\ohm} за да се преполови напонскиот
стрес по компонента и да се подобри температурното совпаѓање на соодносот. Th\'evenin изворната
импеданција претставена на ADC е
\begin{equation}
R_{\mathrm{src}}=R_{\mathrm{top}}\parallel R_{\mathrm{bot}}
   =\frac{110\text{k}\times4.7\text{k}}{110\text{k}+4.7\text{k}}=\SI{4.51}{\kilo\ohm},
\label{eq:rsrc}
\end{equation}
што е клучната големина за прашањето за бафер во Дел~\ref{sec:buffer}.

\subsection{Дали ADC бара op-amp бафер? Квантитативен одговор}\label{sec:buffer}
ADC канал на STM32G4 семплира така што го поврзува изворот, преку внатрешна отпорност на
мултиплексер/прекинувач \(R_{\mathrm{ADC}}\), на кондензатор за семплирање \(C_{\mathrm{ADC}}\) во
текот на програмираното време на семплирање \(t_s\) (Сл.~\ref{fig:adcmodel}). Кондензаторот мора да се
стабилизира во рамки на \(\tfrac12\,\)LSB од изворниот напон. За \(N\)-битна конверзија ова
бара~\cite{an2834,an5346}
\begin{equation}
t_s \;\ge\; \bigl(R_{\mathrm{AIN}}+R_{\mathrm{ADC}}\bigr)\,C_{\mathrm{ADC}}\,\ln\!\bigl(2^{\,N+1}\bigr),
\label{eq:settle}
\end{equation}
каде што \(R_{\mathrm{AIN}}=R_{\mathrm{src}}\) е надворешната изворна импеданција. ST не ги објавува
\(R_{\mathrm{ADC}}\)/\(C_{\mathrm{ADC}}\) директно за G4; наместо тоа техничкиот лист дава табела за
``максимален \(R_{\mathrm{AIN}}\) наспроти време на семплирање'', а упатството на ST е да се избере
времето на семплирање од таа табела за дадената изворна импеданција~\cite{an5346}. Со користење на
моделот на стабилизирање од AN2834 со репрезентативни вредности
\(C_{\mathrm{ADC}}\approx\SI{5}{\pico\farad}\) и \(R_{\mathrm{ADC}}\approx\SI{2}{\kilo\ohm}\), и
\(\ln(2^{13})=9.01\) за 12-битниот режим,
\begin{equation}
t_s \ge (4510+2000)(5\times10^{-12})(9.01)=\SI{293}{\nano\second}.
\end{equation}
При ADC такт од \SI{60}{\mega\hertz} ова е \(17.6\) ADC циклуси, додека STM32G4 нуди поставки за време
на семплирање до \(640.5\) циклуси; самата поставка од \(47.5\) циклуси е \SI{792}{\nano\second} ---
повеќе од \(2.7\times\) од побараното. Бидејќи напоните на сабирницата и фазите бараат само освежување
со kHz-стапка, потрошокот на подолго време на семплирање не чини ништо. Затоа:
\begin{center}
\emph{Не е потребен надворешен op-amp бафер за напонските разделници.}
\end{center}
Изворната импеданција на разделникот се полни до 12-битна точност во рамки на време на семплирање што
е тривијално достапно; мал филтерски кондензатор (Дел~\ref{sec:vfilter}) дополнително ја намалува
ефективната изворна импеданција во моментот на семплирање.

\begin{figure}[t]
\centering
\begin{circuitikz}[scale=1.0,transform shape,american]
\draw (0,1.2) node[anchor=east]{$V_{\mathrm{div}}$} to[R,l={$R_{\mathrm{AIN}}{=}R_{\mathrm{src}}$}] (3,1.2);
\draw (3,1.2) to[R,l=$R_{\mathrm{ADC}}$] (5,1.2);
\draw (5,1.2) to[nos,l=$t_s$] (6.2,1.2);
\draw (6.2,1.2) to[C,l=$C_{\mathrm{ADC}}$] (6.2,0) node[ground]{};
\draw (6.2,1.2) -- (7.0,1.2) node[anchor=west]{SAR};
\end{circuitikz}
\caption{Модел од прв ред на влезниот степен за семплирање на ADC на STM32G4 користен во
\eqref{eq:settle}: изворната импеданција на разделникот $R_{\mathrm{src}}$ на сериска врска со
внатрешната отпорност на прекинувачот го полни кондензаторот за семплирање $C_{\mathrm{ADC}}$ во текот
на прозорецот на семплирање $t_s$.}
\label{fig:adcmodel}
\end{figure}

\begin{table}[t]
\centering
\caption{Op-amp бафер пред ADC: придобивки наспроти трошоци.}
\label{tab:buffer}
\small
\renewcommand{\arraystretch}{1.2}
\begin{tabular}{p{4.0cm} p{4.0cm}}
\toprule
\textbf{Аргументи за бафер} & \textbf{Аргументи против бафер}\\
\midrule
Ниска излезна импеданција $\Rightarrow$ брзо стабилизирање, кратко време на семплирање, висока стапка на мултиплексирано скенирање &
Додава сопствен офсет, дрифт, шум и грешка од ограничен пропусен опсег на \emph{апсолутно} мерење\\
Го изолира разделникот од ``kickback'' поради инјектирање на полнеж во ADC / преслушување меѓу скенираните канали &
Дополнителен трошок на BOM, површина на плочата и декуплирање на напојувањето\\
Овозможува активен anti-alias филтер &
Станува дополнителен товар за калибрација\\
Го спречува резистивното оптоварување на разделникот &
Непотребен тука: бараното $t_s$ е мало, \eqref{eq:settle}\\
\bottomrule
\end{tabular}
\end{table}

Ако некогаш би се посакувал бафер --- на пример за делење на еден разделник меѓу повеќе брзи
мултиплексирани канали, или за додавање активен филтер --- тој треба да се реализира со
\emph{сопствените ресурси на MCU} наместо со надворешен IC: STM32G473 содржи шест операционни
засилувачи што можат да се конфигурираат како follower-и со единечно засилување (или PGA) со нивниот
излез насочен \emph{внатрешно} кон ADC канал, што истовремено го отстранува патот на собирање шум преку
надворешен приклучок и не чини ништо на BOM~\cite{an5306,stm32g473ds}. ST токму оваа follower-кон-ADC
употреба ја документира и препорачува време на семплирање од \SI{200}{\nano\second} при читање на
излезот од внатрешен op-amp~\cite{an5346,an5306}. Компромисот околу баферот е сумиран во
Табела~\ref{tab:buffer}; за тековните напонски канали со ниска стапка заклучокот е да се изостави
баферот и да се користи подолго време на семплирање плюс филтерски кондензатор.

\subsection{Филтрирање на сабирницата и фазите, и заштита}\label{sec:vfilter}
Кондензатор \(C_{\mathrm{f}}\) паралелно со \(R_{\mathrm{bot}}\) формира нископропусен филтер од прв
ред со гранична фреквенција
\begin{equation}
f_c=\frac{1}{2\pi R_{\mathrm{src}}C_{\mathrm{f}}} .
\end{equation}
\textbf{Канал на сабирницата:} \(C_{\mathrm{f}}=\SI{22}{\nano\farad}\) дава
\(f_c=\SI{1.6}{\kilo\hertz}\), што го отстранува склопното бранување (ripple) и шумот од бавното мерење
на сабирницата. \textbf{Фазни канали (back-EMF):} филтерот мора да го отфрли PWM од
\SI{40}{\kilo\hertz} додека ја пропушта основната хармоника на back-EMF. Електричната фреквенција е
\(f_e=\mathrm{RPM}\cdot p/60\) (со \(p\) полови парови); изборот на
\(f_c\approx\)\,2--\SI{3}{\kilo\hertz} го придушува носителот додека ја пропушта основната хармоника,
по цена на фазно доцнење \(\phi=\arctan(f_e/f_c)\). Табела~\ref{tab:phase} го наведува ова доцнење; тоа
е детерминистичко и се компензира во firmware (набљудувачот ја поместува семплираната фаза за \(\phi\)).
За безсензорска комутација базирана на премин низ нула, истото фиксно доцнење се коригира како временски
офсет.

\begin{table}[t]
\centering
\caption{Доцнење на нископропусниот фазен филтер $\phi=\arctan(f_e/f_c)$ (компензирано во firmware).}
\label{tab:phase}
\small
\begin{tabular}{c c c c}
\toprule
$f_c$ & $f_e=\SI{200}{\hertz}$ & $f_e=\SI{500}{\hertz}$ & $f_e=\SI{1000}{\hertz}$\\
\midrule
\SI{2}{\kilo\hertz} & \ang{5.7} & \ang{14.0} & \ang{26.6}\\
\SI{3}{\kilo\hertz} & \ang{3.8} & \ang{9.5}  & \ang{18.4}\\
\bottomrule
\end{tabular}
\end{table}

\textbf{Заштита:} за време на dead-time/спроведување на body-диода фазен јазол се ниша до
\(\approx\SI{-1}{\volt}\), што разделникот го скалира на \(-1/24.4=\SI{-41}{\milli\volt}\) на приклучокот
на ADC --- мало, но Schottky ограничувач (clamp, на пр.\ BAT54S) од јазолот на ADC кон земја и кон
\(V_{\mathrm{DDA}}\) е додаден како осигурување против преодни појави над/под опсегот, при што долниот
отпорник на разделникот ја ограничува струјата на ограничувачот. Истиот ограничувач штити од врвови на
пренапон на каналот на сабирницата. Сл.~\ref{fig:divider} ги прикажува разделникот, филтерот и
ограничувачот.

\begin{figure}[t]
\centering
\begin{circuitikz}[scale=1.05,transform shape,american]
% input and top resistor
\draw (0,4.2) node[anchor=south]{$V_{\mathrm{in}}$ (сабирница / фаза)} -- (0,3.4);
\draw (0,3.4) to[R,l={$R_{\mathrm{top}}$, \SI{110}{\kilo\ohm}}] (0,1.7);
\draw (0,1.7) node[circ]{};
% bottom resistor to ground
\draw (0,1.7) to[R,l={$R_{\mathrm{bot}}$, \SI{4.7}{\kilo\ohm}}] (0,0) node[ground]{};
% horizontal sense line to ADC
\draw (0,1.7) -- (5.2,1.7);
% filter cap tap
\draw (2.3,1.7) node[circ]{} to[C,l=$C_f$] (2.3,0) node[ground]{};
% clamp diodes (ESD pair) at x=3.4
\draw (3.4,1.7) node[circ]{};
\draw (3.4,1.7) to[Do] (3.4,3.4) node[anchor=south]{$V_{\mathrm{DDA}}$};
\draw (3.4,0) node[ground]{} to[Do] (3.4,1.7);
\draw (4.3,3.0) node[anchor=west]{\scriptsize BAT54S};
\draw (5.2,1.7) node[anchor=west]{кон ADC};
\end{circuitikz}
\caption{Напонски разделник за мерење (сооднос $24.4$), филтерски кондензатор од прв ред, и Schottky
ограничувач кон $V_{\mathrm{DDA}}$/GND. Идентични мрежи ги опслужуваат DC-сабирницата
($f_c\approx\SI{1.6}{\kilo\hertz}$) и трите фазни јазли ($f_c\approx$\,2--\SI{3}{\kilo\hertz} за
back-EMF).}
\label{fig:divider}
\end{figure}

\section{RC snubber на склопниот јазол}\label{sec:snubber}

\subsection{Потекло на ѕвонењето (ringing)}
При секое исклучување, паразитната индуктивност \(L_p\) на комутационата јамка (изводи на пакувањето,
PCB траги, ESL на кондензаторот на сабирницата) резонира со излезната капацитивност на MOSFET \(C_o\)
(\(C_{\mathrm{oss}}\) на уредот плюс монтажната капацитивност), произведувајќи придушена осцилација на
склопниот јазол на
\begin{equation}
f_{\mathrm{ring}}=\frac{1}{2\pi\sqrt{L_p C_o}} ,
\label{eq:fring}
\end{equation}
вообичаено во десетици MHz. Ѕвонењето го оптоварува напонот на уредот, зрачи EMI и се спрегнува во
сигналните линии. RC придушувачки snubber поставен паралелно со уредот додава отпорник \(R_s\) (на
сериска врска со кондензатор \(C_s\)) кој ја дисипира резонантната
енергија~\cite{severns,digikey_snubber}.

\subsection{Методологија на дизајн}
Ригорозната постапка, независна од layout-от, го мери ѕвонењето на склопената плоча~\cite{severns,
digikey_snubber}:
\begin{enumerate}[leftmargin=1.4em]
\item Измери го периодот на ѕвонење без snubber \(T_1=1/f_{\mathrm{ring}}\) на склопниот јазол.
\item Додај познат кондензатор \(C_{\mathrm{test}}\) паралелно со уредот и измери го новиот период \(T_2\).
\item Бидејќи \(T_1=2\pi\sqrt{L_pC_o}\) и \(T_2=2\pi\sqrt{L_p(C_o+C_{\mathrm{test}})}\), квадрирањето и
одземањето го елиминира \(L_p\):
\begin{equation}
L_p=\frac{T_2^{2}-T_1^{2}}{4\pi^{2}\,C_{\mathrm{test}}},\qquad
C_o=\frac{T_1^{2}}{4\pi^{2}L_p}.
\label{eq:extract}
\end{equation}
\end{enumerate}
Придушувачкиот отпорник се поставува на карактеристичната импеданција на резонанцата,
\begin{equation}
R_s=\sqrt{\frac{L_p}{C_o}},
\label{eq:rs}
\end{equation}
што дава приближно критично придушување~\cite{digikey_snubber}. Snubber кондензаторот се избира неколку
пати поголем од \(C_o\) така што \(R_s\) да доминира во јамката, а сепак доволно мал за да се ограничи
дисипацијата; моќноста на отпорникот е енергијата складирана во \(C_s\), полнета и празнета еднаш по
склопен циклус,
\begin{equation}
P_{R_s}=C_s\,\Vdcbus^{2}\,\Fswsw .
\label{eq:psnub}
\end{equation}

\subsection{Проценети вредности (да се валидираат на тест-маса)}
Сè уште не постои склопена плоча, па \(L_p\) не може да се измери; вредностите подолу се прва проценка
само за избор на footprint и \emph{мора} да се потврдат со методот со две мерења од \eqref{eq:extract}
на прототипот. Од техничкиот лист на IPT015N10N5, \(C_{\mathrm{oss}}=\SI{1800}{\pico\farad}\) (наведено
при тест-условот од техничкиот лист; \(C_{\mathrm{oss}}\) силно зависи од напонот)~\cite{ipt015ds,ipt015part};
додавањето \(\approx\SI{200}{\pico\farad}\) монтажна капацитивност дава \(C_o\approx\SI{2}{\nano\farad}\).
Земајќи репрезентативно \(L_p\approx\SI{10}{\nano\henry}\) за збиен layout,
\begin{align}
f_{\mathrm{ring}}&=\frac{1}{2\pi\sqrt{(10\text{n})(2\text{n})}}=\SI{35.6}{\mega\hertz},\\
R_s&=\sqrt{\frac{10\text{nH}}{2\text{nF}}}=\SI{2.24}{\ohm}\;\to\;\SI{2.2}{\ohm}\ (\text{E12}).
\end{align}
Со избор на \(C_s=\SI{2.2}{\nano\farad}\) (\(\approx C_o\)), дисипацијата во најлош случај (сабирница од
\SI{60}{\volt}) од \eqref{eq:psnub} е
\begin{equation}
P_{R_s}=2.2\text{n}\times 60^{2}\times 40\text{k}=\SI{0.32}{\watt}\ \text{по snubber},
\end{equation}
т.е.\ \(\le\SI{1.9}{\watt}\) за сите шест прекинувачи ако секој footprint е монтиран (само
\SI{0.08}{\watt} по секој на сабирницата од \SI{30}{\volt}). Се специфицира отпорник со
\(\ge\SI{1}{\watt}\) за резерва.

\subsection{Избор на компоненти}
\(C_s\) мора да биде со C0G/NP0 диелектрик: Class-2 (X7R) керамиките губат поголем дел од својата
капацитивност под DC bias и самозагреваат под високофреквентното бранување, додека C0G е без загуби и
стабилен. Компонента со \SI{2.2}{\nano\farad}, \SI{100}{\volt} C0G е основната (компонента од
\SI{200}{\volt} дава дополнителна резерва против пречекорувањето од ѕвонење). \(R_s\) мора да биде
неиндуктивен: SMD thin/thick-film чип (1206 или 2512), никогаш wirewound, така што неговата сопствена
индуктивност да не го поништи snubber-от~\cite{severns}. RC се поставува директно паралелно со
drain--source на секој MOSFET со најкратка можна јамка. Според спецификацијата, сите шест RC места се
поставени како \textbf{do-not-populate (DNP)} footprint-и: ако прототипот покаже прифатливо ѕвонење тие
остануваат празни, а ако не, се монтираат користејќи ги вредностите измерени на тест-маса од
\eqref{eq:extract}--\eqref{eq:rs}. Сл.~\ref{fig:snub} ги прикажува поставувањето и ефектот.

\begin{figure}[t]
\centering
\begin{circuitikz}[scale=0.95,transform shape,american]
% low-side device as box
\draw (0,3.4) node[anchor=east]{склопен јазол} -- (1.2,3.4);
\draw (1.2,2.6) rectangle (2.2,3.4);
\draw (1.7,3.0) node{\small $Q$};
\draw (1.7,2.6) -- (1.7,1.6) node[ground]{};
\draw (1.7,3.4) -- (1.7,3.4);
% snubber across D-S
\draw (1.7,3.4) -- (3.0,3.4) to[R,l=$R_s$\,\SI{2.2}{\ohm}] (3.0,2.6)
      to[C,l=$C_s$\,\SI{2.2}{\nano\farad}] (3.0,1.6) node[ground]{};
\draw (1.7,3.4) node[circ]{};
\draw (3.0,3.4) node[anchor=south]{\scriptsize DNP RC};
% ring sketch
\begin{scope}[shift={(4.6,1.7)}]
\draw[->] (0,0) -- (3.2,0) node[anchor=north]{\scriptsize $t$};
\draw[->] (0,-0.9) -- (0,1.5) node[anchor=east]{\scriptsize $v$};
\draw[blue,thick] plot coordinates {(0.00,0.85) (0.05,0.73) (0.11,0.46) (0.16,0.10) (0.21,-0.26) (0.27,-0.54) (0.32,-0.68) (0.38,-0.66) (0.43,-0.50) (0.48,-0.24) (0.54,0.06) (0.59,0.33) (0.64,0.50) (0.70,0.56) (0.75,0.49) (0.80,0.31) (0.86,0.08) (0.91,-0.16) (0.96,-0.34) (1.02,-0.44) (1.07,-0.44) (1.12,-0.34) (1.18,-0.17) (1.23,0.03) (1.29,0.20) (1.34,0.33) (1.39,0.37) (1.45,0.33) (1.50,0.22) (1.55,0.06) (1.61,-0.09) (1.66,-0.22) (1.71,-0.29) (1.77,-0.29) (1.82,-0.23) (1.88,-0.12) (1.93,0.01) (1.98,0.13) (2.04,0.21) (2.09,0.24) (2.14,0.22) (2.20,0.15) (2.25,0.05) (2.30,-0.05) (2.36,-0.14) (2.41,-0.19) (2.46,-0.19) (2.52,-0.15) (2.57,-0.08) (2.62,0.00) (2.68,0.08) (2.73,0.14) (2.79,0.16) (2.84,0.15) (2.89,0.10) (2.95,0.04) (3.00,-0.03)};
\draw[red,thick] plot coordinates {(0.00,0.85) (0.05,0.67) (0.11,0.38) (0.16,0.08) (0.21,-0.18) (0.27,-0.35) (0.32,-0.41) (0.38,-0.36) (0.43,-0.25) (0.48,-0.11) (0.54,0.03) (0.59,0.13) (0.64,0.18) (0.70,0.18) (0.75,0.15) (0.80,0.09) (0.86,0.02) (0.91,-0.04) (0.96,-0.07) (1.02,-0.09) (1.07,-0.08) (1.12,-0.06) (1.18,-0.03) (1.23,0.00) (1.29,0.03) (1.34,0.04) (1.39,0.04) (1.45,0.03) (1.50,0.02) (1.55,0.01) (1.61,-0.01) (1.66,-0.02) (1.71,-0.02) (1.77,-0.02) (1.82,-0.01) (1.88,-0.01) (1.93,0.00) (2.04,0.01) (2.20,0.00) (2.62,0.00) (3.00,0.00)};
\node[blue,anchor=west] at (1.3,1.15) {\scriptsize без snubber};
\node[red,anchor=west] at (1.3,0.62) {\scriptsize со snubber};
\end{scope}
\end{circuitikz}
\caption{DNP RC snubber паралелно со секој прекинувач и квалитативниот ефект врз ѕвонењето на
склопниот јазол: придушувачкиот отпорник $R_s=\sqrt{L_p/C_o}$ ја отстранува резонантната енергија.
Прикажаните вредности се први проценки што треба да се заменат со измерени на тест-маса преку
\eqref{eq:extract}.}
\label{fig:snub}
\end{figure}

\section{Стратегија за шум и филтрирање}\label{sec:noise}
Филтрирањето се применува на секоја точка каде што се генерира чувствителен аналоген сигнал или каде
што се инјектира склопен шум:
\begin{itemize}[leftmargin=1.3em]
\item \textbf{Излез на CSA} (Дел~\ref{sec:csafilter}): лесен RC од \SI{723}{\kilo\hertz}
      (\SI{100}{\ohm}/\SI{2.2}{\nano\farad}) по фаза, потпирајќи се на синхроно семплирање наместо на
      anti-alias филтер што внесува доцнење~\cite{an5346}.
\item \textbf{Влез на CSA}: сите влезни сериски отпорници мора да бидат \emph{совпаднати} меѓу IN+ и
      IN-- за да се зачува CMRR; enhanced PWM rejection на INA240 е примарната одбрана против CM
      преодната појава, па тежок влезен филтер намерно се избегнува~\cite{ina240ds}.
\item \textbf{Разделник на сабирницата}: RC од \SI{1.6}{\kilo\hertz} (Дел~\ref{sec:vfilter}).
\item \textbf{Фазни разделници}: RC од 2--\SI{3}{\kilo\hertz} со компензација на фазното доцнење во
      firmware, плюс Schottky ограничувачи.
\item \textbf{Аналогно напојување}: ferrite bead меѓу \(V_{\mathrm{DD}}\) и \(V_{\mathrm{DDA}}\) со
      декуплирање \SI{1}{\micro\farad}\,$+$\,\SI{100}{\nano\farad} на \(V_{\mathrm{DDA}}\); посебно
      \SI{1}{\micro\farad}\,$+$\,\SI{100}{\nano\farad} на \(V_{\mathrm{REF+}}\), и декуплирањето на
      VREFBUF специфицирано од ST ако се користи внатрешната референца~\cite{stm32g473ds}.
\item \textbf{Напојување на INA240}: \SI{100}{\nano\farad}\,$+$\,\SI{1}{\micro\farad} поставени на
      приклучокот.
\item \textbf{Декуплирање на DC-сабирницата}: керамики со ниска ESL паралелно со bulk електролитите,
      блиску до секој half-bridge, за да се минимизира индуктивноста на комутационата јамка \(L_p\) ---
      ова го напаѓа \emph{изворот} на ѕвонењето што snubber-от инаку мора да го придушува~\cite{severns}.
\item \textbf{Заземјување/layout}: единствен (star) спој меѓу аналогната земја (AGND) и енергетската
      земја (PGND); Kelvin мерење на shunt-от; кратки, тесно спрегнати диференцијални сигнални траги; и
      компактна склопна јамка. Овие мерки во layout-от се исто толку важни колку и филтрите со
      компоненти за инвертор со \SI{40}{\kilo\hertz} и \SI{>100}{\ampere}.
\end{itemize}

\section{Спецификација на материјали (BOM, LCSC)}\label{sec:bom}
Табела~\ref{tab:bom} ги наведува компонентите на аналогниот влезен степен со LCSC нарачкини броеви.
Активните уреди (INA240A2PWR, IPT015N10N5) и стандардните чип отпорници/кондензатори се производи со
голем обем што рутински имаат залиха далеку над \num{500} единици; единствената компонента што треба да
се потврди при пуштање на нарачката е shunt-от со голема моќност, за кој е дадена алтернатива составена
од два крајно вообичаени \SI{0.2}{\milli\ohm} \(2512\) елементи паралелно (\(=\SI{0.1}{\milli\ohm}\),
\(\approx\SI{0.93}{\watt}\) по секој). Количините претпоставуваат една трифазна плоча; snubber
компонентите се DNP (само footprint-и).

\begin{table*}[t]
\centering
\caption{Спецификација на материјали за аналогниот влезен степен. Залихата се менува; потврди при пуштање на нарачката.}
\label{tab:bom}
\small
\renewcommand{\arraystretch}{1.2}
\begin{tabular}{p{3.1cm} p{1.0cm} p{3.2cm} p{2.6cm} p{2.0cm} p{2.7cm}}
\toprule
\textbf{Функција} & \textbf{Кол.} & \textbf{Компонента / вредност} & \textbf{Клучен номинал} & \textbf{LCSC} & \textbf{Забелешка}\\
\midrule
Струен засилувач & 3 & TI INA240A2PWR & $G{=}50$, $-4..\SI{80}{\volt}$ CM & C129949 & по еден на фаза~\cite{lcsc_ina}\\
Фазен shunt & 3 & \SI{0.1}{\milli\ohm} метален елемент, \SI{5}{\watt} & $\le\SI{2}{\nano\henry}$ ESL, Kelvin & потврди & на пр.\ Milliohm 6.6$\times$6.6\,mm; види алт.\ подолу~\cite{lcsc_shunt}\\
\;\;(алтернатива за shunt) & 6 & \SI{0.2}{\milli\ohm} \(2512\) $\times2$ \(\parallel\) & \SI{2}{\watt} по сек. & голема залиха & крајно вообичаени; $=\SI{0.1}{\milli\ohm}$\\
Излезен филтер R на CSA & 3 & \SI{100}{\ohm} \(0603\) 1\% & --- & голема залиха & со $C_f$\\
Излезен филтер C на CSA & 3 & \SI{2.2}{\nano\farad} C0G \(0603\) & \SI{50}{\volt} & голема залиха & $f_c\,\SI{723}{\kilo\hertz}$\\
Горен дел на разделник & 4(8) & \SI{110}{\kilo\ohm} (2$\times$\SI{55}{\kilo\ohm}) \(0805\) 1\% & $\ge\SI{150}{\volt}$ & голема залиха & сабирница + 3 фази\\
Долен дел на разделник & 4 & \SI{4.7}{\kilo\ohm} \(0603\) 1\% & --- & голема залиха & сооднос $24.4$\\
Филтер кондензатор на сабирница & 1 & \SI{22}{\nano\farad} C0G/X7R \(0603\) & \SI{50}{\volt} & голема залиха & $f_c\,\SI{1.6}{\kilo\hertz}$\\
Филтер кондензатор на фаза & 3 & \SI{22}{\nano\farad} \(0603\) & \SI{50}{\volt} & голема залиха & $f_c\approx\SI{2.5}{\kilo\hertz}$\\
Ограничувачка диода на ADC & 4 & BAT54S двојна Schottky & \SI{30}{\volt} & голема залиха & кон $V_{\mathrm{DDA}}$/GND\\
Snubber R (DNP) & 6 & \SI{2.2}{\ohm} \(1206\) неинд., $\ge\SI{1}{\watt}$ & --- & голема залиха & само footprint\\
Snubber C (DNP) & 6 & \SI{2.2}{\nano\farad} C0G \(1206\) & \SI{100}{\volt} (\SI{200}{\volt} опц.) & голема залиха & само footprint\\
Ferrite bead на VDDA & 1 & \SI{600}{\ohm}@\SI{100}{\mega\hertz} \(0805\) & --- & голема залиха & $V_{\mathrm{DD}}\!\to\!V_{\mathrm{DDA}}$\\
Сет за декуплирање & --- & \SI{100}{\nano\farad}/\SI{1}{\micro\farad} \(0402/0603\) & \SI{16}{\volt}+ & голема залиха & CSA, VDDA, VREF\\
Енергетски MOSFET (реф.) & 6 & Infineon IPT015N10N5 & \SI{100}{\volt}/\SI{300}{\ampere}/\SI{1.5}{\milli\ohm} & C108964 & енергетски степен~\cite{lcsc_ipt}\\
\bottomrule
\end{tabular}
\end{table*}

\section{Заклучок}
Дизајниран е и трошковно проценет комплетен аналоген влезен степен за \mbox{20--\SI{60}{\volt}},
\mbox{3--\SI{5}{\kilo\watt}}, \SI{40}{\kilo\hertz} BLDC/PMSM FOC инвертор. Континуираната фазна струја
во најлош случај е изведена од опсегот на напојување/моќност како \(\approx\SI{136}{\ampere}\) RMS
(\(\approx\SI{193}{\ampere}\) врв) во работната точка \SI{3}{\kilo\watt}/\SI{20}{\volt}, и искористена
за димензионирање на Kelvin shunt од метален елемент со \SI{0.1}{\milli\ohm}/\SI{5}{\watt}, читан од
INA240A2 ($G{=}50$) со референца на средна скала, давајќи линеарен опсег од \(\pm\SI{290}{\ampere}\)
врв со резолуција од \SI{0.16}{\ampere/LSB} и офсет \(<\SI{0.25}{\ampere}\). In-line топологијата е
оправдана наспроти low-side, high-side, single-shunt и Hall алтернативите со нејзината континуирана,
независна од duty набљудливост, во согласност со TIDA-00913 на TI. За напонските разделници (сооднос
\(24.4\)) \emph{квантитативно} е покажано дека не им треба надворешен бафер: изворот од
\SI{4.5}{\kilo\ohm} се стабилизира до 12-битна точност за \(<\SI{300}{\nano\second}\), далеку во рамки
на достапниот прозорец на семплирање, а шесте внатрешни op-amp засилувачи на STM32G473 остануваат како
резервно решение без дополнителен трошок. RC snubber-от е димензиониран од објавениот
\(C_{\mathrm{oss}}\) и стандардниот метод со две мерења за екстракција на ѕвонењето
(\(R_s=\sqrt{L_p/C_o}\), \(P_{R_s}=C_s\Vdcbus^2\Fswsw\)) и поставен како DNP footprint на сите шест
прекинувачи. Филтрирањето е специфицирано на излезот на CSA, на разделниците, на аналогното
напојување/референца и на DC-сабирницата, заедно со star-ground и Kelvin мерките во layout-от што се
решавачки на ова ниво на моќност. Следниот чекор е пуштање на прототипот: мерење на вистинското ѕвонење
на склопниот јазол за финализирање на snubber-от, верификација на шумот на струјната јамка при синхроно
семплирање, и потврда на залихата на shunt-от при пуштање на BOM.

\begin{thebibliography}{99}
\bibitem{stm32g473ds} STMicroelectronics, ``STM32G473xB/xC/xE datasheet (DS12712),'' ревизија~5, 2023.
[Онлајн]. Достапно на: \url{https://www.st.com/resource/en/datasheet/stm32g473cb.pdf}
\bibitem{ucc27211} Texas Instruments, ``UCC27211A-Q1 120-V boot, 4-A peak half-bridge gate driver,''
datasheet. [Онлајн]. Достапно на: \url{https://www.ti.com/lit/gpn/UCC27211A-Q1}
\bibitem{ipt015ds} Infineon Technologies, ``IPT015N10N5 OptiMOS\texttrademark~5 \SI{100}{\volt}
power transistor,'' технички лист, ревизија~2.4. [Онлајн]. Достапно на:
\url{https://www.infineon.com/assets/row/public/documents/24/49/infineon-ipt015n10n5-datasheet-en.pdf?fileId=5546d4624a75e5f1014ac94680661aff}
\bibitem{ipt015part} Infineon Technologies, ``IPT015N10N5 product page'' (страница на производот)
(\(C_{\mathrm{oss}}=\SI{1800}{\pico\farad}\), \SI{100}{\volt}, \SI{300}{\ampere}). [Онлајн]. Достапно на:
\url{https://www.infineon.com/part/IPT015N10N5}
\bibitem{ina240ds} Texas Instruments, ``INA240 $-4$\,V to $80$\,V, bidirectional, ultra-precise
current-sense amplifier with enhanced PWM rejection (SBOS662),'' технички лист. [Онлајн]. Достапно на:
\url{https://www.ti.com/lit/gpn/INA240}
\bibitem{tida00913} Texas Instruments, ``TIDA-00913: 48-V/10-A three-phase inverter with in-line
shunt-based phase current sensing reference design.'' [Онлајн]. Достапно на:
\url{https://www.ti.com/tool/TIDA-00913}
\bibitem{ti_lowside_threephase} Texas Instruments, ``Low-drift, low-side current measurements for
three-phase systems,'' апликациона белешка (од документацијата на INA240). [Онлајн]. Достапно на:
\url{https://www.ti.com/product/INA240}
\bibitem{an5346} STMicroelectronics, ``AN5346: STM32G4 ADC use tips and recommendations.'' [Онлајн].
Достапно на:
\url{https://www.st.com/resource/en/application_note/an5346-stm32g4-adc-use-tips-and-recommendations-stmicroelectronics.pdf}
\bibitem{an5306} STMicroelectronics, ``AN5306: Operational amplifier (OPAMP) usage in STM32G4 Series.''
[Онлајн]. Достапно на:
\url{https://www.st.com/resource/en/application_note/an5306-operational-amplifier-opamp-usage-in-stm32g4-series-stmicroelectronics.pdf}
\bibitem{an2834} STMicroelectronics, ``AN2834: How to get the best ADC accuracy in STM32
microcontrollers.'' [Онлајн]. Достапно на:
\url{https://www.st.com/resource/en/application_note/an2834-how-to-get-the-best-adc-accuracy-in-stm32-microcontrollers-stmicroelectronics.pdf}
\bibitem{severns} R.~Severns, ``Design of snubbers for power circuits,'' Cornell Dubilier апликациона белешка. [Онлајн]. Достапно на: \url{https://www.cde.com/resources/technical-papers/design.pdf}
\bibitem{digikey_snubber} DigiKey, ``Resistor-capacitor (RC) snubber design for power switches,''
технички напис. [Онлајн]. Достапно на:
\url{https://www.digikey.com/en/articles/resistor-capacitor-rc-snubber-design-for-power-switches}
\bibitem{lcschome} LCSC Electronics, почетна страница на дистрибутер (\num{560000}+ SKU, \(2600+\) производители).
[Онлајн]. Достапно на: \url{https://www.lcsc.com/}
\bibitem{lcsc_ina} LCSC Electronics, ``INA240A2PWR (C129949).'' [Онлајн]. Достапно на:
\url{https://www.lcsc.com/product-detail/C129949.html}
\bibitem{lcsc_ipt} LCSC Electronics, ``IPT015N10N5 (C108964).'' [Онлајн]. Достапно на:
\url{https://www.lcsc.com/product-detail/C108964.html}
\bibitem{lcsc_shunt} LCSC Electronics, ``Категорија: current sense / shunt отпорници.'' [Онлајн]. Достапно на:
\url{https://www.lcsc.com/products/Current-Sense-Resistors-Shunt-Resistors_909.html}
\end{thebibliography}

\end{document}

